\chapter{Conclusion and Future Works}

This project concludes with the successful design, implementation, and evaluation of **FraudNet**, an end-to-end machine learning platform for real-time credit card fraud detection. By combining static feature engineering with sequence-based modeling and local model explainability, the platform provides a robust, low-latency scoring service suitable for financial transaction processing.

\section{Summary of Contributions}

The primary contributions of this project include:
\begin{itemize}
    \item \textbf{Robust Data Engineering}: Implemented a 12-step automated ML pipeline that handles data cleaning, feature engineering, and class imbalance. Stratified, group-aware data splits and training-fold SMOTE oversampling were used to prevent target leakage.
    \item \textbf{Hybrid Ensemble Design}: Developed a custom model blending Random Forest and SimpleRNN component predictions. This hybrid approach captures both static transactional anomalies and temporal sequence dependencies.
    \item \textbf{Low-Latency Scoring API}: Built a FastAPI REST server that scores transaction payloads in under 50 milliseconds.
    \item \textbf{Explainable Web Interface}: Built a Streamlit web application that provides real-time TreeSHAP force plots, transaction simulations, and batch CSV predictions.
\end{itemize}

\section{Limitations}

Despite its strong performance, the current FraudNet system has several limitations:
\begin{itemize}
    \item \textbf{Sequence Length Dependency}: The RNN component relies on a fixed history window of length 6. Transactions with shorter histories cannot leverage sequence predictions, falling back to tabular models.
    \item \textbf{Off-line Model Updates}: The models are trained offline. To adapt to new fraud patterns, the pipeline must retrain the ensemble periodically rather than updating weights in real-time.
    \item \textbf{Explainer Compute Constraints}: Generating SHAP values is computationally expensive. To maintain low latencies, explainability is restricted to the tabular Random Forest component.
\end{itemize}

\section{Future Enhancements}

To address these limitations, future work will focus on:
\begin{itemize}
    \item \textbf{Transformer-Based Sequence Models}: Replacing the SimpleRNN component with self-attention networks (Transformers) to process variable sequence lengths and better capture long-range temporal patterns.
    \item \textbf{Streaming Data Pipeline}: Integrating Apache Kafka or Apache Flink to ingest and process streaming transaction data in real-time, allowing dynamic updates to historical sequence windows.
    \item \textbf{Online Continual Learning}: Implementing semi-supervised or incremental learning algorithms that update model weights continuously as new, verified transaction labels become available.
\end{itemize}
